% Global route: exact part order and a highlighted current stage.
\newcommand{\globalroute}[1]{%
\par\smallskip
\tagstructbegin{tag=Figure,alttext={全书九段路线依次为学习、矩阵、文字、上下文、Transformer、助手、架构、视觉、联想记忆；第八、九段是主线之后的扩展路径。}}
\begin{center}
\begin{tikzpicture}[x=1mm,y=1mm,>=Latex,
 stage/.style={draw=linegray,fill=paper,rounded corners=2pt,minimum width=43mm,minimum height=8mm,font=\scriptsize,align=center}]
\foreach \n/\title in {1/学习,2/矩阵,3/文字,4/上下文,5/Transformer,6/助手,7/架构,8/视觉,9/联想记忆}{
 \pgfmathtruncatemacro{\col}{mod(\n-1,3)}
 \pgfmathtruncatemacro{\row}{int((\n-1)/3)}
 \ifnum\n=#1\relax
   \node[stage,draw=flowcolor,line width=1.1pt,fill=flowcolor!12] (r\n) at (\col*49,-\row*14) {\n\enspace\title\enspace\textbf{你在这里}};
 \else
   \node[stage] (r\n) at (\col*49,-\row*14) {\n\enspace\title};
 \fi
}
\foreach \a/\b in {1/2,2/3,4/5,5/6,7/8,8/9}{\draw[->,flowcolor] (r\a)--(r\b);}
\draw[->,flowcolor] (r3.south)--++(0,-3)--++(-98,0)--(r4.north);
\draw[->,flowcolor] (r6.south)--++(0,-3)--++(-98,0)--(r7.north);
\end{tikzpicture}\par
\tagstructbegin{tag=Caption}{\small\sffamily\color{muted}全书路线 · 沿编号前进；每部分开篇用一行路线标出当前位置}\par\tagstructend
\end{center}
\tagstructend\par\smallskip}

% Chapter numbers and destinations follow the same labels as the main text.
\newcommand{\mapchapter}[2]{\hyperref[intro:#1]{第~\ref*{intro:#1}~章 · #2}}

% One-line route on each part opener; the full route figure is on the reading page.
% The current stage is marked in words, not by colour alone.
\newcommand{\partcrumb}[2]{\ifnum#1=\value{part}\relax{\bfseries\color{blue}#1\,#2（你在这里）}\else#1\,#2\fi}
\newcommand{\partbreadcrumb}{%
\par\smallskip\noindent{\footnotesize\sffamily\color{muted}\textbf{全书路线}\enspace
\partcrumb{1}{学习}\,›\,\partcrumb{2}{矩阵}\,›\,\partcrumb{3}{文字}\,›\,\partcrumb{4}{上下文}\,›\,%
\partcrumb{5}{Transformer}\,›\,\partcrumb{6}{助手}\,›\,\partcrumb{7}{架构}\enspace｜\enspace 扩展：%
\partcrumb{8}{视觉}\,›\,\partcrumb{9}{联想记忆}}\par}

% Part opener: entry only. Review, answers and the hand-off live in \partexit.
% #1 why start here; #2 think-first question; #3 two \item outcomes; #4 map rows.
\newcommand{\partentry}[4]{%
\par\medskip
\begin{tcolorbox}[enhanced,colback=paper,colframe=linegray,boxrule=.5pt,
  arc=2mm,left=5mm,right=5mm,top=3mm,bottom=3mm]
{\sffamily\footnotesize\bfseries\color{muted}为什么从这里开始}\par
#1
\end{tcolorbox}
\begin{tcolorbox}[enhanced,colback=blueback,colframe=linegray,boxrule=.5pt,
  arc=2mm,left=5mm,right=5mm,top=3mm,bottom=3mm]
{\sffamily\bfseries 先想一想}\enspace{\footnotesize\color{muted}凭直觉答就行，本部分末尾会给出答案。}\par
#2
\end{tcolorbox}
\noindent{\sffamily\bfseries\color{blue}学完你能做到}\par\nobreak
\begin{itemize}[itemsep=1pt,topsep=2pt]
#3
\end{itemize}
\par\medskip\noindent{\sffamily\large\bfseries\color{blue}本部分地图}\par\smallskip
\begingroup\small\renewcommand{\arraystretch}{1.65}
\tagpdfsetup{table/header-rows={1}}
\noindent\begin{tabularx}{\linewidth}{@{}>{\raggedright\arraybackslash}p{43mm}>{\raggedright\arraybackslash}X@{}}
\toprule
章节入口 & 本章核心问题\\\midrule
#4
\bottomrule
\end{tabularx}\endgroup
\partbreadcrumb}

% Part exit: answer the opener question, cumulative review, then hand off.
% #1 part key (partq:<key> marks the answer in this part); #2 question;
% #3 one-sentence answer; #4 hand-off to what follows.
\newcommand{\partexit}[4]{%
\par\Needspace{10\baselineskip}\medskip
\begin{tcolorbox}[enhanced,breakable=false,colback=paper,colframe=linegray,boxrule=0pt,
  arc=2mm,left=4mm,right=4mm,top=2mm,bottom=2mm,borderline west={2.5pt}{0pt}{blue},
  fontupper=\small]
{\sffamily\bfseries 回到开篇}\par
\textbf{先想一想：}#2\par
\textbf{答案：}#3\enspace{\footnotesize\hyperref[partq:#1]{§\ref*{partq:#1} · 第~\pageref*{partq:#1}~页}}
\end{tcolorbox}
\csname partreview#1\endcsname
\par\smallskip\noindent{\small\textbf{下一步：}#4}\par}

% ---- 第一部分 ----
\newcommand{\partthinklearning}{三个样本是 $(-1,-1),(0,2),(1,2)$。用 $y=x+1$ 预测，第一项猜高 $1$，第二项猜低 $1$，第三项猜中。前两项的偏差相加为零，就算全猜对了吗？}
\newcommand{\partentrylearning}{\partentry{%
序章把“回答问题”变成了“预测下一个 Token”。可预测的规则从哪里来？现实中的情况多到规则写不尽，所以先退到最小的情形：输入一个数、输出一个数，看样本怎样把公式里的空位填上。
}{\partthinklearning}{%
\item 给一组样本和一条直线，算出“错了多少”，再判断参数该往哪个方向改、改多大一步。
\item 让两层网络从输出的误差倒推每个参数该怎么改；再用没参与训练的样本检查，它是学会了规律，还是只背下了答案。
}{%
\mapchapter{00-functions}{函数与损失} & 参数怎样改变预测，损失怎样比较整体表现？\\
\mapchapter{16a-gradient-descent}{梯度下降} & 沿哪个方向、走多大一步才能降低损失？\\
\mapchapter{07-functions-relu}{非线性} & 为什么两层直线仍不够，ReLU 改变了什么？\\
\mapchapter{17-backpropagation}{反向传播} & 怎样把输出误差逐层传回每个参数？\\
\mapchapter{18-generalization}{泛化} & 为什么训练成绩好，还要看验证集？\\
}}
\newcommand{\partexitlearning}{\partexit{learning}{\partthinklearning}{%
前两个样本仍然猜错。正负残差会相抵；把残差平方后取平均，得到均方误差 $2/3$，才把两项偏差都计入了分数。
}{%
同样的乘加会排进向量和矩阵，一次算完多个特征、多个样本。
}}

% ---- 第二部分 ----
\newcommand{\partthinkmatrices}{三套房各有面积和房龄两个数，估价公式完全相同。要逐套算三次，还是能把三套房摆成一张表，一次算完？}
\newcommand{\partentrymatrices}{\partentry{%
上一部分的函数一次只吃一个数。可估一套房子的价格，要同时看面积和房龄；实际任务还要一次估一批房子。
}{\partthinkmatrices}{%
\item 不算具体数值，先判断两个矩阵能否相乘、结果是几行几列。
\item 把一批样本写成矩阵，手算一次两层网络，并说明各行共用同一组权重，为什么仍得到不同的输出。
}{%
\mapchapter{01-vectors}{向量与点积} & 怎样按需求给文章打分？为什么放大数字会改变排名，归一化又改变了什么？\\
\mapchapter{03-matrix-multiplication}{矩阵乘法} & 一行乘一列，怎样组成一整张结果表？\\
\mapchapter{10-mlp}{矩阵 MLP} & 怎样让一批输入共享两层权重和激活函数？\\
}}
\newcommand{\partexitmatrices}{\partexit{matrices}{\partthinkmatrices}{%
能一次算完。三套房摆成三行的输入矩阵 $X$，乘同一个权重矩阵 $W$，$XW$ 的每一行就是一套房的结果。
}{%
矩阵里装的都是数。下一部分让文字也成为矩阵中的一行。
}}

% ---- 第三部分 ----
\newcommand{\partthinktokens}{给“春天、花、开了”分别编号 $2$、$4$、$5$。能说“开了”是“春天”的 $2.5$ 倍吗？编号能直接当作特征来计算吗？}
\newcommand{\partentrytokens}{\partentry{%
矩阵只认数字，模型要读的却是“春天 花 开了”这样的文字。
}{\partthinktokens}{%
\item 把一句话切成 Token、写出编号，并为每个位置写出要预测的下一个 Token。
\item 把一行分数变成概率并算出损失，分清一个 Token 的编号、向量和输出概率是三种不同的表示。
}{%
\mapchapter{20a-tokenizer}{Token 与 ID} & “春天 花 开了”怎样切分并编号？\\
\mapchapter{20b-task-contract}{预测任务} & 三个输入位置分别要预测哪一个下一 Token？\\
\mapchapter{20c-classification-loss}{分类损失} & 一行分数怎样变成概率与损失？主线五个候选又怎样计分？\\
\mapchapter{20-embeddings}{Embedding} & 为什么用 ID 查向量，而不直接计算 ID 的距离？\\
}}
\newcommand{\partexittokens}{\partexit{tokens}{\partthinktokens}{%
不能。编号只表示身份，大小没有意义；直接当特征会引入虚假的倍数关系。所以要按编号查出一行可训练的向量。
}{%
现在每个 Token 是矩阵中的一行，但每一行还只看得到自己。下一部分让各行交换信息。
}}

% ---- 第四部分 ----
\newcommand{\partthinkcontext}{“春天，花，开了”和“冬天，窗，开了”的最后一个词相同。只看“开了”这一个词，能判断后面该接什么吗？}
\newcommand{\partentrycontext}{\partentry{%
每个 Token 已经是一行向量，可模型逐行计算：处理“开了”这一行时，它看不到前面的词。
}{\partthinkcontext}{%
\item 用一组对照例子说明，只看当前词为什么不够。
\item 说明怎样让一个位置按需读取前文，以及为什么只能读已经出现的词。
}{%
\mapchapter{21a-context-gap}{上下文缺口与 RNN} & 逐行 MLP 为什么看不到别的 Token？怎样把历史压进逐步传递的状态？\\
\mapchapter{21c-naive-attention}{Attention 与 Q/K/V} & 直接读取怎样计算，三组投影怎样改变匹配与内容？\\
\mapchapter{11-softmax-memory}{批量与因果注意力} & 所有位置怎样一起算？因果 Mask 怎样只读前缀？\\
\mapchapter{22a-position-multihead}{位置与多头} & 怎样标明顺序，并行使用不同匹配方式？\\
}}
\newcommand{\partexitcontext}{\partexit{context}{\partthinkcontext}{%
不能。两句的当前词相同，同一个函数只看这一行，输出必然相同；要区分两句，必须读到前文。
}{%
Attention 只是一个读取步骤。下一部分把它和 FFN、残差、输出头组装成能训练、能生成的完整模型。
}}

% ---- 第五部分 ----
\newcommand{\partthinktransformer}{“开了”这一行读完上下文，得到的仍是一行数字。词表里有五个候选，怎样从这一行数字得到每个候选的概率？}
\newcommand{\partentrytransformer}{\partentry{%
Attention 已能让“开了”读到“春天、花”。但读完只得到一行数字，还不是对下一个词的预测，也还不会从错误中学习。
}{\partthinktransformer}{%
\item 沿“春天、花、开了”走完一次前向，得到概率和损失，并指出训练更新了哪些共用的参数。
\item 说明生成时怎样把新词接回输入，以及为什么不必每一步都从头重算。
}{%
\mapchapter{22-transformer-architecture}{完整结构} & Attention、Norm 与残差怎样连接？\\
\mapchapter{10b-ffn-output-head}{FFN 与输出头} & 特征变换和词表打分为何要分开？\\
\mapchapter{22b-token-to-output}{三维前向} & 三个 Token 怎样得到三行五类概率？\\
\mapchapter{22c-language-training}{联合训练} & 三个位置怎样共同更新共享参数？\\
\mapchapter{23-llm-systems}{逐词生成} & 下一 Token 如何成为下一轮的输入？\\
\mapchapter{23a-from-toy-to-gpt}{真实尺度} & 教学网络与 GPT 的参数、词表和层数差在哪？\\
}}
\newcommand{\partexittransformer}{\partexit{transformer}{\partthinktransformer}{%
先经 FFN 逐行变换，再由输出头 $W_{\rm out}$ 给五个候选各打一个分数（logits），最后用 Softmax 把分数变成概率。
}{%
现在的模型会续写。下一部分让它接上资料、工具和授权，成为能完成任务的助手。
}}

% ---- 第六部分 ----
\newcommand{\partthinkassistant}{让模型写季度报告，它流畅地报出一个第三季度收入，可它从没见过这份财报。这个数字可信吗？该怎么核实？}
\newcommand{\partentryassistant}{\partentry{%
能续写，不等于能完成任务：报告里的数字要有出处，对外的动作要经过授权。
}{\partthinkassistant}{%
\item 区分临时给样例、再训练模型和提供检索资料，各自改变了系统的哪一部分。
\item 给报告中的每个数字找到来源，并说明工具调用由谁提出、由谁执行、失败了怎么办。
}{%
\mapchapter{23aa-pretraining-icl}{ICL} & 更新权重和临时提供样例有何区别？\\
\mapchapter{23c-post-training}{训练反馈与推理时计算} & 奖励怎样更新参数，验证又怎样筛选本次候选？\\
\mapchapter{23d-agent-systems}{RAG} & 报告中的数字怎样找到可核查的来源？\\
\mapchapter{23d-agent-tools}{工具与 Agent} & 谁提出调用，谁执行，怎样处理失败？\\
\mapchapter{23d-agent-protocols}{协作边界} & 协议、Skill 和多个 Agent 怎样分工？\\
\mapchapter{23b-deployment-and-ecosystem}{数据、训练与部署} & 语料如何进入训练，训练和上线各花什么资源？\\
}}
\newcommand{\partexitassistant}{\partexit{assistant}{\partthinkassistant}{%
不可信。流畅只说明续写的概率高，不说明数字有出处；要先检索真实报表、逐项核对来源，再做计算。
}{%
系统层面接上了证据与工具。下一部分回到模型内部，看现代架构为什么不断替换部件。
}}

% ---- 第七部分 ----
\newcommand{\partthinkarchitecture}{生成时，模型为每个旧 Token 缓存一份 K 和 V（\hyperref[intro:23-llm-systems]{第~\ref*{intro:23-llm-systems}~章}的 KV Cache）。有 $4$ 个提问头，能不能让它们共用 $2$ 份 K/V？会省下什么，又可能损失什么？}
\newcommand{\partentryarchitecture}{\partentry{%
“忒修斯之船”问：船的部件逐块换掉以后，还是原来那条船吗？今天的大模型几乎把位置表示、归一化、Attention 和 FFN 都换过一轮。是否还算同一种架构，要看计算结构；换得值不值，要看相同条件下的质量、延迟与成本。
}{\partthinkarchitecture}{%
\item 指出现代大模型替换了原版 Transformer 的哪些部件，各自为了解决什么瓶颈。
\item 遇到“新架构更快”的说法时，列出比较时必须保持一致的条件。
}{%
\mapchapter{22d-modern-transformers}{现代架构} & 首读 RoPE、RMSNorm、GQA，分别改了哪一步？其余专题按需选读。\\
\mapchapter{23e-tech-judgement}{技术判断} & 怎样用同条件实验判断一项改动是否值得？\\
}}
\newcommand{\partexitarchitecture}{\partexit{architecture}{\partthinkarchitecture}{%
能，这就是 GQA：每两个提问头共用一份 K/V，缓存减半。它是“完全独立”与“全部共享”之间的折中，质量损失多少，要在同条件下对照。
}{%
主线到此结束。后面是两条扩展路径：\hyperref[part:vision]{视觉部分}改造模型的输入，\hyperref[part:memory]{联想记忆部分}换一个角度重看 Attention；也可以直接进入附录。
}}

% ---- 第八部分 ----
\newcommand{\partthinkvision}{一张 $30\times30$ 的灰度图展平成 $900$ 个数，送进全连接层。哪些信息变得不明显了？}
\newcommand{\partentryvision}{\partentry{%
这是一条扩展路径，前面的助手与架构两部分不依赖它。Transformer 读的是一串向量，图片却是一张二维网格。
}{\partthinkvision}{%
\item 说明小窗口扫描图片时，为什么不同位置能共用一组参数。
\item 把一张 $4\times4$ 小图切成四块，写出每一步的形状，将视觉向量作为前缀，与文字向量拼接。
}{%
\mapchapter{19-cnn}{CNN} & 小卷积核怎样扫描图像并共享参数？\\
\mapchapter{19b-vit}{ViT 与视觉桥接} & 图像块怎样成为 Token，再接到语言模型？\\
}}
\newcommand{\partexitvision}{\partexit{vision}{\partthinkvision}{%
相邻像素的位置关系不再显式保留；同一个边缘出现在不同位置，全连接层还要分别学习。CNN 用小窗口扫描并共享参数，ViT 把图切块后当作 Token。
}{%
下一部分换一个角度，把 Attention 看成“先存入、再取出”的记忆；也可以直接进入附录。
}}

% ---- 第九部分 ----
\newcommand{\partthinkmemory}{往同一张表里写入两条记录，再用第一条的钥匙去读。读出的只有第一条的内容，还是会混进第二条？}
\newcommand{\partentrymemory}{\partentry{%
这是第二条扩展路径，主线各章不依赖它。前面把 Attention 当作“按需读取前文”；这里换一个角度，把它看成“先存入、再取出”的记忆。
}{\partthinkmemory}{%
\item 把两条记录写进同一张表，算出读取时它们互相干扰了多少。
\item 解释按读出的误差修正写入时，为什么修正一条会扰动其他记录。
}{%
\mapchapter{04-key-value-outer-product}{写入与误差修正} & 怎样写入、读出串扰，再联立求解或逐条修正？\\
\mapchapter{09-nonlinearity-and-memory}{非线性检索与容量} & 筛选怎样改变读出，{\symboliconfont ★}随机容量估计有哪些条件？\\
}}
\newcommand{\partexitmemory}{\partexit{memory}{\partthinkmemory}{%
两把钥匙不垂直、匹配分数不为零时，会混进第二条，这就是串扰；两把钥匙垂直时互不干扰。
}{%
正文到此结束。终章回到开篇的四个 Token；附录提供公式、参数、题库、答案与来源。
}}
